\documentclass[11pt,oneside]{article}
\usepackage[T1]{fontenc}
\usepackage[utf8]{inputenc}
\usepackage{lmodern}
\usepackage[a4paper,margin=27mm,headheight=15pt]{geometry}
\usepackage{microtype,amsmath,amssymb,amsthm,mathtools,aliascnt}
\usepackage{booktabs,array,longtable,enumitem}
\usepackage[dvipsnames]{xcolor}
\usepackage{fancyhdr}
\usepackage{hyperref}
\hypersetup{colorlinks=true,linkcolor=MidnightBlue,citecolor=MidnightBlue,
 urlcolor=MidnightBlue,pdftitle={Omnific-Preserving Automorphisms},
 pdfauthor={OpenAI ChatGPT},pdfsubject={Convex-scale stabilizers, definable constants, and nondefinable monomials}}
\usepackage[nameinlink,noabbrev]{cleveref}
\pagestyle{fancy}\fancyhf{}
\fancyhead[L]{\small Omnific-preserving automorphisms}
\fancyhead[R]{\small\thepage}
\renewcommand{\headrulewidth}{0.3pt}
\setlength{\parskip}{3pt}\setlength{\parindent}{1.2em}
\setlength{\emergencystretch}{2em}
\setcounter{tocdepth}{2}
\numberwithin{equation}{section}
\newtheorem{theorem}{Theorem}[section]
\theoremstyle{plain}
\newaliascnt{proposition}{theorem}
\newtheorem{proposition}[proposition]{Proposition}
\aliascntresetthe{proposition}
\crefname{proposition}{proposition}{propositions}
\Crefname{proposition}{Proposition}{Propositions}
\newaliascnt{lemma}{theorem}
\newtheorem{lemma}[lemma]{Lemma}
\aliascntresetthe{lemma}
\crefname{lemma}{lemma}{lemmas}
\Crefname{lemma}{Lemma}{Lemmas}
\newaliascnt{corollary}{theorem}
\newtheorem{corollary}[corollary]{Corollary}
\aliascntresetthe{corollary}
\crefname{corollary}{corollary}{corollarys}
\Crefname{corollary}{Corollary}{Corollarys}
\theoremstyle{definition}
\newaliascnt{definition}{theorem}
\newtheorem{definition}[definition]{Definition}
\aliascntresetthe{definition}
\crefname{definition}{definition}{definitions}
\Crefname{definition}{Definition}{Definitions}
\newaliascnt{example}{theorem}
\newtheorem{example}[example]{Example}
\aliascntresetthe{example}
\crefname{example}{example}{examples}
\Crefname{example}{Example}{Examples}
\theoremstyle{remark}
\newaliascnt{remark}{theorem}
\newtheorem{remark}[remark]{Remark}
\aliascntresetthe{remark}
\crefname{remark}{remark}{remarks}
\Crefname{remark}{Remark}{Remarks}
\newcommand{\No}{\mathbf{No}}
\newcommand{\Oz}{\mathbf{Oz}}
\newcommand{\R}{\mathbb R}
\newcommand{\C}{\mathbb C}
\newcommand{\Q}{\mathbb Q}
\newcommand{\Z}{\mathbb Z}
\newcommand{\N}{\mathbb N}
\newcommand{\J}{\mathcal J}
\newcommand{\OO}{\mathcal O}
\newcommand{\mm}{\mathfrak m}
\newcommand{\LL}{\mathfrak L}
\newcommand{\UU}{\mathcal U}
\newcommand{\FF}{\mathfrak F}
\newcommand{\Aut}{\operatorname{Aut}}
\newcommand{\Hom}{\operatorname{Hom}}
\newcommand{\supp}{\operatorname{supp}}
\newcommand{\lc}{\operatorname{lc}}
\newcommand{\ct}{\operatorname{ct}}
\newcommand{\id}{\operatorname{id}}
\newcommand{\Fix}{\operatorname{Fix}}
\newcommand{\Exp}{\operatorname{Exp}}
\newcommand{\Log}{\operatorname{Log}}
\newcommand{\BCH}{\operatorname{BCH}}
\newcommand{\ord}{\operatorname{ord}}
\newcommand{\bb}{\mathbin{\ll}}
\newcommand{\pin}{\texttt{fb5c4b530e3ec04b5661347d51d5382526d5aa02}}
\newcommand{\status}[1]{\par\smallskip\noindent\textit{Provenance and scope.} #1\par\smallskip}

\begin{document}
\begin{titlepage}
\centering
\vspace*{13mm}
{\Large A research article for the Surreal project\par}
\vspace{12mm}
{\huge\bfseries Omnific-Preserving\\[3mm] Automorphisms\par}
\vspace{8mm}
{\Large Convex-scale stabilizers, definable constants,\\[2mm]
 and nondefinable monomials\par}
\vspace{11mm}
{\large September 23, 2026\par}
\vspace{8mm}
\begin{minipage}{0.92\textwidth}
\begin{abstract}
We investigate which coefficient-fixing, strongly additive,
leading-term-preserving automorphisms
of a Hahn field preserve its omnific-type integer part. For a divisible
ordered exponent group $\Gamma$, write $C_\delta$ for the principal convex
subgroup generated by $\delta>0$. Our main classification says that a
contracting strong derivation preserves the nonpositive-support ring exactly
when the coefficient functional attached to each positive shift $\delta$
vanishes on $C_\delta$. The established exponential--logarithm correspondence
then gives an exact description of the integer-part stabilizer.

This criterion yields a rank dichotomy and a sharp solvability theorem:
the leading-term-preserving stabilizer is trivial precisely in
Archimedean rank one; in finite ordered rank $r$ its derived length is at
most $r-1$, with equality for the reverse-lexicographic group $\Q^r$.
For the full surreal field, the stabilizer of every set of parameters is
nonsolvable and contains an explicitly constructed copy of the additive
group $\R^\kappa$ for every set cardinal $\kappa$.

A separate definability argument, building on omnific Pell rigidity,
recovers $\R$, the purely infinite ideal, and truncation at exponent zero
from the pair $(\No,\Oz)$. Nevertheless, no set of parameters makes the
Conway monomials definable, and no definable rule can choose one
representative of each natural-valuation class. The common fixed field of
the strong omnific-preserving $1$-automorphisms is exactly $\R$.
Real-coefficient constructions extend to surcomplex numbers and preserve
both conjugation and the Gaussian omnific integers.
\end{abstract}
\end{minipage}
\vfill
\begin{minipage}{0.92\textwidth}
\small
Prepared with OpenAI ChatGPT. Full mathematical arguments are supplied,
with imported theorems and proposed contributions distinguished. This is
an unrefereed research draft, not a claim of Lean verification or exhaustive
historical priority. The accompanying exact-arithmetic program checks
finite identities only. Repository comparison is pinned to the revision
recorded in Appendix~\ref{app:audit}.
\end{minipage}
\end{titlepage}
\pagenumbering{roman}
{\small\setlength{\parskip}{0pt}\tableofcontents}
\clearpage
\pagenumbering{arabic}

\section{The question and the contribution}\label{sec:intro}

A Conway normal form distinguishes three kinds of information: coefficients,
exponents, and the monomial representatives of those exponents. The omnific
integers retain the purely infinite part and an integer constant coefficient.
It is natural to ask how much of the normal-form presentation is determined
by the ambient field together with this distinguished ring.

The answer developed here is asymmetric. The coefficient field and the
cut between negative and nonnegative Hahn exponents are recoverable. The
monomial cross-section is not. In fact, even an arbitrarily large
\emph{set} of named surreal parameters leaves enough automorphisms to move
monomials and to defeat every definable choice of valuation representatives.
The mechanism is an explicit family of derivations that change a higher
Archimedean scale by a lower one without ever moving a purely infinite term
across the constant-term boundary.

\subsection{Main statements in a common notation}

Let $k$ have characteristic zero and let $\Gamma\ne0$ be a divisible ordered
abelian group. Use the Hahn convention
\[
 K=k((t^\Gamma)),\qquad v(f)=\min\supp(f),\qquad
 B=k((t^{\Gamma_{\le0}})),\qquad
 \J=k((t^{\Gamma_{<0}})).
\]
For a unital subring $D\subseteq k$, put $A_D=D\oplus\J$.
For the surreal field, $k=\R$, $\Gamma=\No$, and
$t^\gamma=\omega^{-\gamma}$, so $A_\Z=\Oz$.
A $1$-automorphism fixes every leading term, not merely every valuation.
The adjective \emph{strong} means that every Hahn-summable family remains
summable and that its sum is preserved.

The principal results are as follows.
\begin{enumerate}[label=\textup{(\arabic*)},leftmargin=*,itemsep=4pt]
\item A contracting strong $k$-derivation with monomial expansion
\[
 \partial(t^\gamma)
 =t^\gamma\sum_{\delta>0}c_\delta(\gamma)t^\delta
\]
satisfies $\partial(B)\subseteq B$ exactly when
$c_\delta(C_\delta)=0$ for every $\delta>0$. This is also exactly the
condition that $\Exp(\partial)$ preserve $A_D$
(\cref{thm:stabilizer}).
\item The resulting group $\UU_{A_D}$ is trivial exactly when $\Gamma$
is Archimedean. Its finite-rank derived-length bound $r-1$ is optimal
(\cref{thm:dichotomy,thm:rank,thm:sharp}).
\item In the first-order pair $(\No,+,\cdot,<,\Oz)$, the classes $\J$,
$B$, $\R$, and $\Z$, the natural valuation ring, and zero-cut truncation
are parameter-free definable (\cref{thm:recovery}).
\item For every set $P\subseteq\No$, a strong $1$-automorphism preserving
$\Oz$ fixes $P$ pointwise while moving a Conway monomial to a nonmonomial.
No $P$-definable selection of natural-valuation representatives exists
(\cref{thm:parameters,thm:nosection}).
\item The common fixed field of $\UU_\Oz$ is $\R$. Its pointwise
stabilizer of every set $P$ contains $\R^\kappa$ for every set cardinal
$\kappa$, and is nonsolvable (\cref{thm:fixed,thm:product,thm:nonsolvable}).
\end{enumerate}

These results concern genuinely infinite normal forms and the full
integer-part predicate, not merely polynomial substitutions in one infinite
variable. The distinction is essential: a translation may preserve
$k[t^{-1}]$ and still fail to preserve $B$ after divisible exponents are
allowed.

\subsection{What is imported, and what is proposed here?}

Conway normal forms and the omnific integer part are classical
\cite{Conway,Gonshor}. The ambient theory of strong Hahn-field automorphisms
is developed by Kuhlmann--Serra \cite{KS}. We import the
exponential--logarithm correspondence for contracting strong derivations
from Bagayoko--Krapp--Kuhlmann--Panazzolo--Serra
\cite[Theorem 3.13]{BKKPS}; its surreal extension is stated by
Kaplan--Krapp--Serra \cite[Fact 4.2]{KKS}. Those correspondence and
factorization theorems are \emph{not} contributions of this article.

The inspected repository already has a report on surcomplex automorphisms
and an omnific Diophantine report proving constant-product and Pell rigidity
\cite{RepoAut,RepoOz}. We explicitly reuse those ideas. The additional
contributions proposed here are the exact convex-subgroup stabilizer
criterion, its sharp derived-length consequences, the definability recovery
via Pell divisibility and idealizers, and the parameter-uniform
nondefinability and largeness theorems. These precise combined statements
were not located in the inspected primary sources or relevant repository
material. That is a bounded comparison, not a proof of historical priority.

No named longstanding conjecture is declared solved here. Instead, we give
new candidate theorems with complete arguments, as permitted by the research
brief. Their value rests on the proofs and on their distinction from the
imported general machinery, rather than on an unsupported claim of novelty.

\subsection{Class size and the meaning of an automorphism collection}

Unless the full surreal specialization is explicitly being discussed,
exponent groups are set-sized. This gives ordinary fields and ordinary
groups. The abstract stabilizer proof applies to $\No$ using the cited
class version of the operator correspondence. When
$\Gamma=\No$, we work in a class-theoretic background such as NBG with global
choice. Every individual Hahn support is a set. An automorphism is an
individual class map; we do not form a set, or an NBG class-object, whose
members are all such proper-class graphs. A statement about
$\UU_\Oz$ is interpreted as a statement about individual maps, their
compositions, or a set-indexed action encoded by one class relation.

Likewise, a common-fixed-field statement is pointwise: an element is fixed
by every automorphism under discussion exactly when it has the stated
property. Nonsolvability means that no ordinary finite derived length
annihilates all iterated commutators. For valuation representatives we use
an equivalence relation on nonzero field elements, not a collection of
proper-class cosets treated as sets. These conventions are important for
\cref{thm:nosection,thm:product}.

\section{Hahn support, the boundary ring, and strong maps}\label{sec:prelim}

\subsection{Summability and the sign convention}

An element of $K=k((t^\Gamma))$ is a formal sum
\[
 f=\sum_{\gamma\in S}a_\gamma t^\gamma
\]
with $S\subseteq\Gamma$ well ordered. A family $(f_i)_{i\in I}$ is
\emph{Hahn summable} if the union of its supports is well ordered and each
exponent occurs in only finitely many of the supports. Its sum is then
formed coefficientwise using only finite sums in $k$. Throughout,
$v(0)=\infty$.

We use the standard support lemma: if $S$ is well ordered and
$T\subseteq\Gamma_{>0}$ is well ordered, then
$S+\langle T\rangle_{\N}$ is well ordered; each prescribed exponent
receives only finitely many contributions from a choice of an element of
$S$ and a finite word in $T$. Here the empty word is allowed. This is the
Neumann support principle, in the word form needed for evaluations of
operator power series; see the summability treatment in \cite{BKKPS}.
No ordinary convergence of real or complex coefficient series is being
substituted for this finiteness assertion.

For actual surreals we always take
\begin{equation}\label{eq:sign}
 t^\gamma=\omega^{-\gamma}.
\end{equation}
Thus negative $t$-exponents are purely infinite terms, and positive
$t$-exponents are infinitesimal terms. In particular,
$\J$ is supported at \emph{negative}, not positive, $t$-exponents.

\subsection{Four different rings or additive pieces}

The decompositions relevant to this article are
\begin{equation}\label{eq:decomp}
 B=k\oplus\J,\qquad A_D=D\oplus\J,\qquad
 \OO=k\oplus\mm,
\end{equation}
where
\[
 \OO=\{f:v(f)\ge0\},\qquad \mm=\{f:v(f)>0\}.
\]
The displayed direct sums are additive decompositions. The set $\J$ is a
nonunital ideal of $B$ and of $A_D$, and $\mm$ is the maximal ideal of the
valuation ring $\OO$. By contrast, $B$ is generally not a valuation ring.
The map $\ct:B\to k$ is a ring retraction and restricts to $A_D\to D$.
Constant-coefficient extraction on the \emph{whole} field is not a ring
homomorphism.

For later reference define the additive projections
\[
 P_-(f)=\sum_{\gamma<0}a_\gamma t^\gamma,
 \qquad P_0(f)=a_0,
 \qquad P_+(f)=\sum_{\gamma>0}a_\gamma t^\gamma.
\]
Every subset of a well-ordered support is well ordered, so these are
legitimate Hahn series.

\begin{lemma}[Constant-product rigidity]\label{lem:units}
If $f,g\in B$ and $fg\in k^\times$, then $f,g\in k^\times$.
Consequently $B^\times=k^\times$.
\end{lemma}
\begin{proof}
Both $v(f)$ and $v(g)$ are nonpositive and their sum is zero. Thus both
are zero. A nonempty support contained in $\Gamma_{\le0}$ with least
member zero consists only of zero. Hence both series are constant.
\end{proof}
\status{This elementary rigidity argument is already used in the repository's
omnific Diophantine report \cite{RepoOz}; we reproduce it because the
model-theoretic application below needs its exact scope.}

\begin{definition}[Principal convex scale]
For $\delta>0$, let
\begin{equation}\label{eq:convex}
 C_\delta=\{\gamma\in\Gamma: |\gamma|\le n\delta
                 \text{ for some }n\in\N\}.
\end{equation}
Write $\delta\bb\epsilon$ when $n\delta<\epsilon$ for every positive
ordinary integer $n$. Thus $C_\delta\subsetneq C_\epsilon$ exactly when
$\delta\bb\epsilon$. The group $C_\delta$ is divisible because
$\Gamma$ is divisible.
\end{definition}

\subsection{The imported operator correspondence}

A $k$-linear map is \emph{strong} if it preserves Hahn summability and
commutes with the sums of all summable families. A derivation $\partial$
is \emph{contracting} if
\[
 v(\partial f)>v(f)\qquad(f\ne0).
\]
An automorphism $\sigma$ is a \emph{$1$-automorphism} if
$v(\sigma f-f)>v(f)$ for every nonzero $f$. Equivalently, it preserves
both $v(f)$ and $\lc(f)$.

We use the following established result in exactly this form.
\begin{theorem}[Strong exponential--logarithm correspondence,
\cite{BKKPS,KKS}]\label{thm:import}
For a characteristic-zero Hahn field, contracting strong $k$-derivations
correspond bijectively to strong $k$-linear $1$-automorphisms by
\begin{align}
 \Exp(\partial)(f)&=\sum_{n\ge0}\frac{\partial^n(f)}{n!},\label{eq:exp}\\
 \Log(\sigma)(f)&=\sum_{n\ge1}\frac{(-1)^{n+1}}n(\sigma-\id)^n(f).
 \label{eq:log}
\end{align}
The families are strongly summable, and composition corresponds to the
Baker--Campbell--Hausdorff operation on derivations. The surreal version
uses set-sized supports and individual class operators.
\end{theorem}

We write $\Exp$ for an \emph{operator} exponential. The notation
$\exp(u)=\sum u^n/n!$ for $v(u)>0$ denotes an ordinary Hahn-series
exponential. Neither notation, by itself, asserts preservation of the
Gonshor exponential on the whole surreal field.

\section{The exact convex-scale stabilizer}\label{sec:stabilizer}

\subsection{Coefficient functionals of a derivation}

For a contracting strong derivation $\partial$ define
\begin{equation}\label{eq:coeff}
 c_\delta(\gamma)=[t^{\gamma+\delta}]\partial(t^\gamma),
 \qquad \delta>0.
\end{equation}
This is a family of scalar functions indexed by positive shifts.
For each fixed $\gamma$, the nonzero terms have a set-sized well-ordered
support. No assumption that all possible shifts form one globally
well-ordered set is made.

\begin{lemma}\label{lem:additive}
Each $c_\delta:\Gamma\to k$ is additive, hence $\Q$-linear, and
\begin{equation}\label{eq:homogeneous}
 \partial(t^\gamma)=t^\gamma\sum_{\delta>0}c_\delta(\gamma)t^\delta.
\end{equation}
If $c_\delta$ vanishes on $(-\delta,0)$, it vanishes on $C_\delta$.
\end{lemma}
\begin{proof}
The product rule gives
\[
 t^{-\gamma-\eta}\partial(t^{\gamma+\eta})
 =t^{-\gamma}\partial(t^\gamma)+t^{-\eta}\partial(t^\eta).
\]
Taking coefficient $\delta$ proves additivity. Characteristic zero and
divisibility imply $c_\delta(q\gamma)=q c_\delta(\gamma)$.

For $h\in C_\delta$, choose a positive integer $N$ such that
$|h|/N<\delta$. Divisibility provides $h/N\in\Gamma$.
If $h\ne0$, one of $h/N$ and $-h/N$ lies in $(-\delta,0)$.
Additivity and multiplication by $N$ imply $c_\delta(h)=0$.
The case $h=0$ is immediate.
\end{proof}

\begin{theorem}[Boundary-preserving derivations]\label{thm:derivation}
For a contracting strong $k$-derivation $\partial$ the following are
equivalent:
\begin{equation}\label{eq:criterion}
 \partial(B)\subseteq B
 \quad\Longleftrightarrow\quad
 c_\delta|_{C_\delta}=0\quad\hbox{for every }\delta>0.
\end{equation}
Under these conditions $\partial(B)\subseteq\J$,
$\partial(\J)\subseteq\J$, and $\partial$ commutes with all three
projections $P_-,P_0,P_+$.
\end{theorem}
\begin{proof}
Suppose first that $\partial(B)\subseteq B$. For
$-\delta<\gamma<0$, the element $t^\gamma$ lies in $B$, whereas
$\gamma+\delta>0$. The coefficient of $t^{\gamma+\delta}$ in its
derivative must therefore vanish. Hence $c_\delta$ vanishes on
$(-\delta,0)$, and \cref{lem:additive} gives its vanishing on
$C_\delta$.

Conversely, suppose the coefficient condition holds. If $\gamma<0$ and
$c_\delta(\gamma)\ne0$, then $\gamma\notin C_\delta$. Therefore
$-\gamma>n\delta$ for every positive integer $n$, and in particular
\begin{equation}\label{eq:nocross}
 \gamma+\delta<0.
\end{equation}
Thus every nonzero term of $\partial(t^\gamma)$ remains strictly on the
negative side. The derivative of a constant is zero. Strongness lets us
sum over every negative support, proving $\partial(B)\subseteq\J$.
Positive exponents remain positive because the shifts are positive.
The image of the negative, zero, and positive pieces consequently lies,
respectively, in the negative piece, zero, and the positive piece.
This proves all projection identities.
\end{proof}

The condition excludes a possible boundary term, not only a term that
crosses the boundary. In particular $c_\delta(-\delta)=0$ follows by
testing $-\delta/2$ and using additivity. This small point is essential
when the constant coefficients of the integer part lie in $D$ rather
than all of $k$.

\begin{theorem}[Exact integer-part stabilizer]\label{thm:stabilizer}
Let $\sigma$ be a strong $k$-linear $1$-automorphism of $K$, and let
$\partial=\Log(\sigma)$. The following conditions are equivalent:
\begin{enumerate}[label=\textup{(\roman*)},leftmargin=*,itemsep=2pt]
\item $\sigma(A_D)=A_D$;
\item $\sigma(B)=B$;
\item $\partial(B)\subseteq B$;
\item $c_\delta|_{C_\delta}=0$ for every $\delta>0$;
\item $\sigma P_-=P_-\sigma$.
\end{enumerate}
When they hold, $\sigma$ fixes constant coefficients and preserves $\J$
and $\mm$ setwise. The subgroup $\UU_{A_D}$ is independent of the choice
of $D\subseteq k$ and is exactly $\Exp(\LL_A)$, where
\[
 \LL_A=\{\partial:\partial\text{ is a contracting strong }k
 \text{-derivation and }c_\delta(C_\delta)=0\text{ for all }\delta>0\}.
\]
\end{theorem}
\begin{proof}
Assume (i). For every $\gamma<0$, the monomial $t^\gamma$ lies in $A_D$,
so its image lies in $B$. Strong $k$-linearity then implies
$\sigma(B)\subseteq B$. Applying the same argument to $\sigma^{-1}$
proves equality, giving (ii).

Under (ii), the operator $T=\sigma-\id$ sends $B$ to $B$, hence so
does each $T^n$. The logarithm in \eqref{eq:log} is a strongly summable
family. A summable family of series supported in $\Gamma_{\le0}$ has
its sum supported there. Thus $\partial(B)\subseteq B$, proving (iii).
Notice that this step is taken in the $k$-algebra $B$, not directly in
$A_D$; there is no unjustified division by the integers $n$ in $D$.

The equivalence of (iii) and (iv), together with preservation of the three
support pieces, is \cref{thm:derivation}. If (iv) holds, every iterate
of $\partial$ preserves $\J$ and annihilates $k$. Hence both
$\Exp(\partial)$ and $\Exp(-\partial)$ preserve $D\oplus\J$ and
commute with $P_-$. This gives (i) and (v).

Finally, (v) implies that $\sigma$ preserves the range of $P_-$,
namely $\J$, and the same identity for the inverse gives equality.
Since $\sigma$ fixes $k$, it preserves $A_D$. Constant coefficients are
unchanged because the negative and positive pieces remain separate and
constants are fixed.
\end{proof}

\begin{remark}[What has, and has not, been classified]
The theorem is an exact classification inside the \emph{strong,
coefficient-fixing, leading-term-preserving} automorphisms. It does not
assert that every abstract field automorphism preserving $A_D$ is strong.
The coefficient functionals in \eqref{eq:coeff} come from an existing
strong derivation; arbitrary functions satisfying the vanishing condition
cannot be assembled into an operator without also checking summability.
\end{remark}

\section{Explicit shears and a rank-one obstruction}\label{sec:shears}

\subsection{One positive shift}

Let $\theta:\Gamma\to(k,+)$ be additive and let $\delta>0$. Define
\begin{equation}\label{eq:singleD}
 D_{\delta,\theta}\left(\sum_\gamma a_\gamma t^\gamma\right)
 =\sum_\gamma a_\gamma\theta(\gamma)t^{\gamma+\delta}.
\end{equation}
Translation of a well-ordered support remains well ordered. The same
observation for a summable family proves that this is a strong map.
Additivity of $\theta$ proves the product rule on monomials, and strong
summability extends it to products of arbitrary Hahn series. Thus
$D_{\delta,\theta}$ is a contracting strong derivation.

\begin{proposition}[Homogeneous shear]\label{prop:shear}
Suppose $\theta(C_\delta)=0$. Then
\begin{align}
 D_{\delta,\theta}^{\,n}(t^\gamma)
   &=\theta(\gamma)^n t^{\gamma+n\delta},\label{eq:iterate}\\
 \sigma_{\delta,\theta}(t^\gamma)
   &:={\Exp(D_{\delta,\theta})}(t^\gamma)
     =t^\gamma\exp\bigl(\theta(\gamma)t^\delta\bigr).
     \label{eq:shear}
\end{align}
The strong extension of \eqref{eq:shear} is an element of $\UU_{A_D}$,
with inverse $\sigma_{\delta,-\theta}$. For fixed $\delta$ these shears
satisfy
\[
 \sigma_{\delta,\theta}\sigma_{\delta,\psi}
 =\sigma_{\delta,\theta+\psi}
 \qquad\bigl(\theta(C_\delta)=\psi(C_\delta)=0\bigr).
\]
Conversely, $\Exp(D_{\delta,\theta})$ preserves $A_D$ only if
$\theta(C_\delta)=0$.
\end{proposition}
\begin{proof}
The hypothesis gives $\theta(\delta)=0$, hence
$\theta(\gamma+n\delta)=\theta(\gamma)$, proving
\eqref{eq:iterate} by induction. Formula \eqref{eq:shear} follows.
The inverse and addition laws follow either from commuting derivations
or directly from the fact that both shears fix $t^\delta$.
The stabilizer assertion and its converse follow from
\cref{thm:stabilizer}, since the only shift functional is $\theta$.
\end{proof}

For clarity, the summability behind the displayed formula can be checked
without importing the general operator theorem. For a series supported
on $S$, the terms of its image are supported in $S+\N\delta$. The
Neumann support principle gives well-ordering and only finitely many
contributions to each coefficient, including across all $n$. This also
justifies termwise products, the inverse, and preservation of arbitrary
summable families. For $\gamma<0$ with $\theta(\gamma)\ne0$, every
$\gamma+n\delta$ remains negative because $-\gamma>n\delta$ for all
ordinary $n$.

\subsection{Uniformly supported families of shifts}

\begin{proposition}[A direct construction criterion]\label{prop:uniform}
Let $S\subseteq\Gamma_{>0}$ be a set-sized well-ordered set, and let
$(\theta_s)_{s\in S}$ be additive maps $\Gamma\to k$. Then
\begin{equation}\label{eq:uniform}
 D(f)=\sum_{\gamma\in\supp(f)}\sum_{s\in S}
       a_\gamma\theta_s(\gamma)t^{\gamma+s}
\end{equation}
is a contracting strong derivation. If $\theta_s(C_s)=0$ for every
$s$, its exponential lies in $\UU_{A_D}$.
The family of all terms used in $\Exp(D)(f)$ is supported in
$\supp(f)+\langle S\rangle_\N$ and is Hahn summable.
\end{proposition}
\begin{proof}
The support lemma applies first to $\supp(f)+S$ and then to all finite
words in $S$. Every coefficient in \eqref{eq:uniform}, in its iterates,
and in products of such expressions is a finite sum. Additivity of each
$\theta_s$ proves the product rule. Replacing a single input support by
the union of the supports of a summable input family proves strongness:
that union is well ordered, and each original exponent has only finitely
many input occurrences. Each nonzero output exponent is larger than its
input exponent, proving contraction. The final assertion follows either
from \cref{thm:stabilizer} or from the no-crossing argument applied at
each step of a word in $S$.
\end{proof}

This proposition will supply the entire parameter-fixing construction on
$\No$ below. All shifts used there form one set-sized well-ordered set,
so those examples do not require an appeal to a proper-class version of
an unrestricted operator-evaluation theorem.

\subsection{Why ordinary translations are not enough}

\begin{example}[Preserving polynomials is weaker than preserving the integer part]
Take $\Gamma=\Q$, $\delta=1$, and $\theta(\gamma)=\gamma$.
Then
\[
 D^n(t^\gamma)
 =\gamma(\gamma+1)\cdots(\gamma+n-1)t^{\gamma+n},
\]
so
\begin{equation}\label{eq:badtranslation}
 \Exp(D)(t^\gamma)=t^\gamma(1-t)^{-\gamma}.
\end{equation}
For $\gamma=-1$ this sends $t^{-1}$ to $t^{-1}-1$ and preserves the
polynomial ring $k[t^{-1}]$. But
\[
 \Exp(D)(t^{-1/2})
 =t^{-1/2}(1-t)^{1/2}
 =t^{-1/2}-\tfrac12t^{1/2}-\tfrac18t^{3/2}-\cdots
\]
has positive support. Hence it does not preserve $B$ or $A_D$.
The failure is exactly $\theta(C_1)\ne0$.
\end{example}

Divisibility of $\Gamma$ is essential to the theorem, not a cosmetic
convenience. With $\Gamma=\Z$, the nonpositive-support ring is
$k[t^{-1}]$, and the translation just exhibited does preserve it.
The proof of \cref{lem:additive} cannot be used in $\Z$: there is no
exponent $-1/2$ at which to detect the boundary crossing.

\begin{theorem}[Archimedean dichotomy]\label{thm:dichotomy}
For a nonzero divisible ordered abelian group $\Gamma$,
\[
 \UU_{A_D}=\{\id\}
 \quad\Longleftrightarrow\quad
 \Gamma\text{ is Archimedean}.
\]
\end{theorem}
\begin{proof}
If $\Gamma$ is Archimedean, $C_\delta=\Gamma$ for every $\delta>0$.
The stabilizer criterion forces every coefficient functional to vanish.
Thus the logarithm is zero and the automorphism is the identity.

If $\Gamma$ is not Archimedean, some $C_\delta$ is proper. The quotient
$\Gamma/C_\delta$ is a nonzero $\Q$-vector space. Choose a nonzero
$\Q$-linear map from this quotient to $\Q$, and compose with
$\Q\subseteq k$. The resulting nonzero $\theta$ vanishes on
$C_\delta$, so \cref{prop:shear} gives a nonidentity stabilizer element.
For set-sized $\Gamma$ the choice of a vector-space functional is an
ordinary application of a basis; explicit surreal coefficient functionals
will replace it in the proper-class case.
\end{proof}

\section{The rest of the strong valued stabilizer}\label{sec:factorization}

The leading-term restriction is useful but is not the only natural
notion of automorphism. Say that $\sigma$ is \emph{valued} if it induces
an order automorphism $\tau$ of $\Gamma$ through
\[
 v(\sigma f)=\tau(v(f))\qquad(f\ne0).
\]
We retain strongness and pointwise fixing of $k$.

For a character $\chi:(\Gamma,+)\to(k^\times,\cdot)$ and an ordered
group automorphism $\tau$, define
\begin{equation}\label{eq:monomialauto}
 M_{\chi,\tau}\left(\sum a_\gamma t^\gamma\right)
   =\sum a_\gamma\chi(\gamma)t^{\tau(\gamma)}.
\end{equation}
Order preservation gives admissible supports; the character identity
gives multiplicativity. This is a strong field automorphism, and it
preserves $A_D$ because $\tau$ preserves the signs of exponents and
$\chi(0)=1$.

\begin{theorem}[Strong valued factorization relative to the integer part]
\label{thm:factorization}
Every strong valued $k$-fixing automorphism $\sigma$ with
$\sigma(A_D)=A_D$ has a unique factorization
\[
 \sigma=M_{\chi,\tau}\circ u,\qquad u\in\UU_{A_D}.
\]
Here $\tau$ is its induced exponent-group automorphism and
$\chi(\gamma)=\lc(\sigma(t^\gamma))$.
\end{theorem}
\begin{proof}
Multiplication of monomials and multiplicativity of leading coefficients
show that $\chi(\gamma+\eta)=\chi(\gamma)\chi(\eta)$.
The map $u=M_{\chi,\tau}^{-1}\sigma$ fixes each monomial's leading
term. For a general series, the valuations of the images of its monomial
terms are strictly ordered by $\tau$. Strongness therefore shows that
the least exponent and leading coefficient of $u(f)$ agree with those
of $f$. Thus $u$ is a $1$-automorphism. Both factors preserve $A_D$,
so $u\in\UU_{A_D}$. The definitions of $\tau$ and $\chi$ show uniqueness.
\end{proof}
\status{The ambient Hahn-field factorization is established in \cite{KS}.
The proof above records its elementary restriction to the particular
integer-part stabilizer. The new information is the explicit description
of $u$ in \cref{thm:stabilizer}, not the general factorization pattern.}

For $k=\R$ and divisible $\Gamma$, every character has positive values:
$\chi(\gamma)=\chi(\gamma/2)^2>0$. Consequently the monomial maps
and the $1$-automorphisms preserve the Hahn ordering. In the case
$\Gamma=\Q$, the factor $u$ is absent, and every such real
integer-part-preserving automorphism has
\[
 t^q\longmapsto c^q t^{\lambda q},
 \qquad c>0,\quad\lambda\in\Q_{>0}.
\]
Here $\chi(q)=c^q$ follows from uniqueness of positive roots, and ordered
automorphisms of the additive group $\Q$ are positive rational scalings.
This is a concrete sense in which higher-rank shears are new phenomena,
not just complicated one-variable substitutions.

\section{Convex rank and sharp solvability}\label{sec:rank}

\subsection{The bracket law}

\begin{lemma}[Homogeneous commutators]\label{lem:bracket}
For additive $\theta,\psi:\Gamma\to k$ and positive $\delta,\epsilon$,
\begin{equation}\label{eq:bracket}
 [D_{\delta,\theta},D_{\epsilon,\psi}]
 =D_{\delta+\epsilon,\,\theta(\epsilon)\psi-\psi(\delta)\theta}.
\end{equation}
If $\theta(C_\delta)=0$ and $\psi(C_\epsilon)=0$, the bracket is zero
when $C_\delta=C_\epsilon$. If $C_\delta\subsetneq C_\epsilon$, it
has shift in the $C_\epsilon$ scale and its functional vanishes on
$C_\epsilon$.
\end{lemma}
\begin{proof}
On $t^\gamma$, the coefficient of $t^{\gamma+\delta+\epsilon}$ in
the two compositions differs by
\[
 \psi(\gamma)\theta(\gamma+\epsilon)
 -\theta(\gamma)\psi(\gamma+\delta)
 =\theta(\epsilon)\psi(\gamma)-\psi(\delta)\theta(\gamma).
\]
This proves \eqref{eq:bracket}. If the convex subgroups agree, each
functional kills the other's shift. If $C_\delta\subsetneq C_\epsilon$,
then $\psi(\delta)=0$ and $C_{\delta+\epsilon}=C_\epsilon$.
The remaining functional is a scalar multiple of $\psi$.
\end{proof}

The same computation applies to general strong derivations after expansion
into their shift coefficients. The coefficient at a prescribed output is
well defined: the two compositions are strongly summable, and the
intermediate product of the two relative monomial derivatives is a product
of two Hahn series. Its finite coefficient sums justify splitting the
Leibniz terms in \eqref{eq:bracket}. No unrestricted double sum of
scalar coefficients is being assumed.

\subsection{A filtration by convex scales}

We say that $\Gamma$ has \emph{finite ordered rank $r$} when its distinct
nonzero principal convex subgroups are precisely
\[
 0<C_1<C_2<\cdots<C_r=\Gamma.
\]
For a positive shift $\delta$, its layer is the index $j$ such that
$C_\delta=C_j$. A shift in layer $r$ contributes no nonzero derivation
in $\LL_A$, because its coefficient functional must vanish on all of
$\Gamma$.

For $1\le j\le r$, let $\FF_j$ consist of the derivations in $\LL_A$
whose nonzero shifts have layers at least $j$. Then
\begin{equation}\label{eq:filtration}
 \LL_A=\FF_1\supseteq\FF_2\supseteq\cdots\supseteq\FF_r=0.
\end{equation}
These are Lie ideals, closed under the strongly summable operator sums
that occur in evaluations. To see the support assertion, decompose the
images of monomials into their individual terms and retain only terms in
the prescribed shift layers. Selecting a subfamily preserves summability.
The new coefficient functions remain additive, so the selected operators
are again strong derivations.

\begin{lemma}\label{lem:rankfiltration}
The filtration satisfies
\[
 [\FF_j,\FF_j]\subseteq\FF_{j+1},\qquad
 [\LL_A,\FF_j]\subseteq\FF_j.
\]
\end{lemma}
\begin{proof}
Apply \cref{lem:bracket} to shift pairs. Equal layers cancel. Unequal
layers produce a term in the larger layer. If both input layers are at
least $j$, a nonzero result must therefore have layer at least $j+1$.
If just one layer is at least $j$, the larger layer is at least $j$.
The support justification above permits summation over all terms.
\end{proof}

For a Lie algebra set $\LL^{(0)}=\LL$ and
$\LL^{(n+1)}=[\LL^{(n)},\LL^{(n)}]$, where brackets generate a
$k$-linear subspace. For a group use the ordinary derived subgroups.
The identity group and zero Lie algebra have derived length zero.

\begin{theorem}[Finite-rank bound]\label{thm:rank}
If $\Gamma$ has finite ordered rank $r$, then both $\LL_A$ and
$\UU_{A_D}$ have derived length at most $r-1$.
\end{theorem}
\begin{proof}
The Lie assertion follows inductively from
$\LL_A^{(n)}\subseteq\FF_{n+1}$.
For the group, put $G_j=\Exp(\FF_j)$. The BCH theorem makes each
$G_j$ a subgroup. The logarithm of a group commutator of two elements of
$G_j$ is a strongly summable Lie series all of whose terms contain a
bracket of two elements of $\FF_j$. By \cref{lem:rankfiltration},
these terms lie in $\FF_{j+1}$, and their sum does too. Consequently
$[G_j,G_j]\subseteq G_{j+1}$. Since $G_r=\{\id\}$, the asserted
bound follows.
\end{proof}

\subsection{Sharpness, including the abstract group}

Let $\Gamma=\Q^r$ with the ordering in which the last nonzero coordinate
of a vector decides its sign. For its coordinate vectors,
\[
 0<e_1\bb e_2\bb\cdots\bb e_r.
\]
Let $\theta_i$ extract coordinate $i$. For $j<i$ set
$X_{j,i}=D_{e_j,\theta_i}$. These derivations satisfy the stabilizer
criterion, since $C_{e_j}$ is spanned by $e_1,\ldots,e_j$.
Define
\begin{equation}\label{eq:ak}
 a_0=0,\qquad a_k=\sum_{h=1}^k2^{k-h}e_h,
 \qquad Y_{k;j,i}=D_{a_k+e_j,\theta_i}\quad(k<j<i\le r).
\end{equation}

\begin{lemma}\label{lem:balanced}
For every allowed triple, $Y_{k;j,i}\in\LL_A^{(k)}$, and
\begin{equation}\label{eq:balanced}
 [Y_{k-1;k,j},Y_{k-1;j,i}]=Y_{k;j,i}\qquad(k\ge1).
\end{equation}
\end{lemma}
\begin{proof}
At $k=0$ this is the definition of $X_{j,i}$. In the bracket in
\eqref{eq:balanced}, $\theta_j(a_{k-1}+e_j)=1$ whereas
$\theta_i(a_{k-1}+e_k)=0$. Its shift is
$2a_{k-1}+e_k+e_j=a_k+e_j$. The bracket law and induction prove the claim.
\end{proof}

The passage from a nonzero iterated Lie bracket to a nonzero group
commutator deserves a justification: the abstract derived subgroup is
not silently identified with the exponential of the algebraic derived
Lie subalgebra.

\begin{lemma}[Mixed-parameter detection]\label{lem:mixed}
Take finitely many homogeneous positive-shift derivations, and form an
iterated Lie bracket $W$ of them. Replace every leaf occurrence by
$\Exp(s_\ell D_\ell)$, with independent scalar parameters, and every
bracket by the group word $ghg^{-1}h^{-1}$. Call the resulting word
$\mathcal W(\mathbf s)$. In its action on a monomial, the coefficient
of the product of all leaf parameters is the action of $W$.
If $W\ne0$, some scalar specialization $\mathbf s\in k^m$ gives a
nonidentity group word.
\end{lemma}
\begin{proof}
Work first with formal parameters. The identity
\[
 \begin{aligned}
 &(1+sD+\cdots)(1+uE+\cdots)(1-sD+\cdots)(1-uE+\cdots)\\
 &\hspace{12mm}=1+su[D,E]+\text{terms of other multidegrees}
 \end{aligned}
\]
proves the first assertion by induction on the bracket tree. The two
subtrees have disjoint parameter sets; their group words become the
identity when any required leaf parameter is zero. Hence the coefficient
of the squarefree product cannot receive contributions from a term
omitting a leaf or repeating its parameter.

Choose a monomial and an output exponent on which $W$ acts nontrivially.
Only a finite set $S$ of positive shifts occurs at the leaves. The
Neumann support principle makes the coefficient at that output a
\emph{polynomial} in the parameters: only finitely many words in $S$
can reach it. By the first assertion this polynomial is nonzero.
An infinite field admits a scalar tuple at which a nonzero polynomial
does not vanish. The specialized group word is therefore nonidentity.
This argument also works for repeated leaves because their occurrences
were assigned separate parameters.
\end{proof}

\begin{theorem}[Sharp derived length]\label{thm:sharp}
For the reverse-lexicographic group $\Gamma=\Q^r$, over every
characteristic-zero field $k$ and every unital $D\subseteq k$, the group
$\UU_{A_D}$ and its Lie algebra have derived length exactly $r-1$.
\end{theorem}
\begin{proof}
For $r=1$ use \cref{thm:dichotomy}. If $r\ge2$,
$Y_{r-2;r-1,r}\ne0$ belongs to $\LL_A^{(r-2)}$ by
\cref{lem:balanced}. Thus the upper bound in \cref{thm:rank} is sharp
for the Lie algebra. Its balanced bracket expression gives, by
\cref{lem:mixed}, a nontrivial group word in the $(r-2)$nd derived
subgroup. This proves sharpness for the group as well.
\end{proof}

\begin{example}[Metabelian need not mean nilpotent]\label{ex:nilpotent}
In rank three put $X=D_{e_1,\theta_2}$ and $Y=D_{e_2,\theta_3}$.
Repeated use of \eqref{eq:bracket} gives
\[
 (\operatorname{ad}X)^n(Y)=D_{e_2+ne_1,\theta_3}\ne0
 \qquad(n\ge0).
\]
Thus the stabilizer Lie algebra is metabelian but not nilpotent.
Applying \cref{lem:mixed} to these brackets shows that the group is
also not nilpotent. Rank two gives a nontrivial abelian group; rank
three already separates solvability from nilpotence.
\end{example}

\section{Recovering the real constants from the omnific predicate}\label{sec:definability}

We now specialize to $K=\No$, $A=\Oz$, and $k=\R$.
All definability assertions in this section are first-order in the pair
\[
 \mathcal S=(\No,+,\cdot,0,1,<,A).
\]
They do not refer to birthday, simplicity, the omega-map, a derivation,
or a named predicate for the ordinary real numbers. In fact the order
symbol can be omitted, since positivity in a real closed field is
definable by nonzero squares.

\subsection{A Pell-divisibility definition of the purely infinite ideal}

\begin{lemma}[Omnific Pell rigidity]\label{lem:pell}
If $u,v\in A$ satisfy $u^2-2v^2=1$, then $u,v\in\Z$.
There are ordinary solutions with positive $u$ larger than any prescribed
ordinary integer.
\end{lemma}
\begin{proof}
Both $u+\sqrt2v$ and $u-\sqrt2v$ lie in $B$, and their product is one.
By \cref{lem:units} both are real constants. Hence $u,v\in\R\cap A=\Z$.
For unbounded ordinary solutions, start with $(u_0,v_0)=(1,0)$ and use
\[
 u_{m+1}=3u_m+4v_m,\qquad v_{m+1}=2u_m+3v_m.
\]
A direct calculation preserves $u_m^2-2v_m^2=1$, and
$u_m\ge3^m$ for $m\ge0$. This establishes the required unboundedness
without using an additional Diophantine theorem.
\end{proof}
\status{The rigidity statement is already in \cite{RepoOz}. The following
use of all Pell first coordinates to define an ideal is the additional
construction here.}

Define a formula $\Phi_J(a)$ in the pair by
\begin{equation}\label{eq:Jformula}
 \Phi_J(a)\ \Longleftrightarrow\ A(a)\ \wedge\
 \forall u,v\in A\;\Bigl[u^2-2v^2=1\ \Longrightarrow\
                  \exists b\in A:\ a=ub\Bigr].
\end{equation}
Restricted quantifiers abbreviate the corresponding implications or
conjunctions with the predicate $A$. This is one fixed first-order formula,
not an infinite conjunction indexed by ordinary integers.

\begin{proposition}\label{prop:Jdef}
For every $a\in\No$, the formula $\Phi_J(a)$ holds exactly when
$a\in\J$.
\end{proposition}
\begin{proof}
If $a\in\J$, every Pell first coordinate $u$ is a nonzero ordinary
integer by \cref{lem:pell}. Then $a/u\in\J\subseteq A$, providing
$b$. Conversely let $a\in A$ satisfy the formula and write
$a=n+j$ with $n\in\Z$, $j\in\J$. The equality $a=ub$ implies,
after taking constant coefficients, $n=u\ct(b)$ in $\Z$.
If $n\ne0$, choose a positive Pell first coordinate $u>|n|$.
It cannot divide $n$, a contradiction. Thus $n=0$ and $a\in\J$.
\end{proof}

\subsection{The idealizer and its units}

\begin{lemma}[Recovering the boundary ring]\label{lem:idealizer}
Within $\No$,
\begin{equation}\label{eq:idealizer}
 B=(\J:\J):=\{x:\forall j\in\J,\ xj\in\J\}.
\end{equation}
\end{lemma}
\begin{proof}
An element of $B$ has only nonpositive exponents, so multiplying a
strictly negative-support series keeps the support strictly negative.
This proves one inclusion. If $x\notin B$, its support contains some
$\eta>0$. Set $j=t^{-\eta/2}\in\J$. Multiplication by this monomial
translates every exponent injectively. In particular $xj$ has the
nonzero coefficient formerly at $\eta$ now at $\eta/2>0$, so it is
not in $\J$. This proves the reverse inclusion.
\end{proof}

\begin{theorem}[Parameter-free recovery]\label{thm:recovery}
In the pair $\mathcal S$, all of the following are parameter-free
definable: $\J$, $B$, $\R$, $\Z$, the natural valuation ring $\OO$,
its maximal ideal $\mm$, and the functions $P_-$, $P_0$, $P_+$.
In particular,
\begin{equation}\label{eq:Rformula}
 \R=\{0\}\cup B^\times,
 \qquad \Z=A\cap\R.
\end{equation}
\end{theorem}
\begin{proof}
The definition of $\J$ is \eqref{eq:Jformula}, and that of $B$ is
\eqref{eq:idealizer}. The condition of being a unit of $B$ is
\[
 B(x)\wedge\exists y\,(B(y)\wedge xy=1).
\]
By \cref{lem:units}, the units are exactly $\R^\times$, proving
\eqref{eq:Rformula}.

Using the now defined real constants, the valuation predicates are
\begin{align}
 x\in\OO&\ \Longleftrightarrow\
     \exists r\in\R\ (r>0\ \wedge\ -r\le x\le r),\label{eq:Oformula}\\
 x\in\mm&\ \Longleftrightarrow\
     \forall r\in\R\ (r>0\ \Longrightarrow\ -r<x<r).
     \label{eq:mformula}
\end{align}
The normal-form ordering proves that these are respectively the
nonnegative- and positive-support pieces.

For each $x$, its negative-support part is the unique $j$ satisfying
\begin{equation}\label{eq:Pformula}
 j\in\J\quad\hbox{and}\quad x-j\in\OO.
\end{equation}
Existence follows by splitting its support at zero; uniqueness follows
from $\J\cap\OO=\{0\}$. After finding $j$, the constant coefficient
is the unique $c\in\R$ such that $x-j-c\in\mm$, since
$\R\cap\mm=\{0\}$. These formulas define the graphs of $P_-$ and
$P_0$. Finally $P_+(x)=x-P_-(x)-P_0(x)$.
\end{proof}

\begin{corollary}[Set-sized Hahn-field version]
The parameter-free recovery in \cref{thm:recovery} holds also in the
ordinary set-sized structure
$\bigl(\R((t^\Gamma)),\Z\oplus\J\bigr)$ for every nonzero divisible
ordered abelian group $\Gamma$.
\end{corollary}
\begin{proof}
The proof uses only real coefficients, the Hahn ordering, divisibility
of exponents in the idealizer argument, and the ordinary Pell recurrence.
It does not require any proper-class bounding principle.
\end{proof}

\begin{corollary}[Consequences for all pair automorphisms]\label{cor:allautos}
Every field automorphism of $\No$ preserving $\Oz$ setwise fixes the
ordinary reals pointwise, preserves the natural valuation ring, and
commutes with the three zero-cut projections. This holds without a
strongness hypothesis. Among strong automorphisms, therefore,
\cref{thm:factorization} describes the entire $\Oz$-stabilizer; no
additional coefficient-fixing or valued assumption is needed.
\end{corollary}
\begin{proof}
Definability gives setwise preservation of $\R$, $\OO$, and the
projection graphs. Every field automorphism of a real closed field
preserves order, because positive elements are precisely the nonzero
squares. Its restriction to $\R$ is then an order-preserving field
automorphism fixing $\Q$, hence fixes each real by its rational cut.
The definable valuation ring determines the natural-valuation
equivalence and its ordered value group. For a strong automorphism the
hypotheses of \cref{thm:factorization} follow.
\end{proof}

\begin{remark}[A definable coefficient is not a global ring retraction]
The function $P_0$ is definable on all of $\No$, but
$P_0(t^{-1})P_0(t)=0$ while $P_0(t^{-1}t)=1$.
Its multiplicativity on $A$ and $B$ is a separate support fact.
Likewise, definability of $P_-$ does not yet define the set of all
monomials. The next section gives an obstruction to doing so.
\end{remark}

\begin{corollary}[Truncation at a named monomial]
For every monomial $m=t^\alpha$, truncation below $\alpha$ is definable
with the single parameter $m$, by
\[
 T_{<\alpha}(f)=m\,P_-(f/m).
\]
\end{corollary}
\begin{proof}
Dividing by $m$ translates an exponent $\gamma$ to $\gamma-\alpha$;
$P_-$ keeps exactly the terms with $\gamma<\alpha$. Multiplication by
$m$ restores their exponents.
\end{proof}

This corollary concerns a specified boundary supplied as a monomial
parameter. It does not provide a uniform definable choice of such a
parameter at every valuation. The distinction matters when comparing
with work on expansions by truncation, such as \cite{Camacho}.

\section{Automorphisms invisible to any set of parameters}\label{sec:parameters}

\subsection{Coefficient functionals on the exponent group}

For $\eta\in\No$, let
\begin{equation}\label{eq:qeta}
 q_\eta(\gamma)=[\omega^\eta]\gamma\qquad(\gamma\in\No).
\end{equation}
The coefficient here belongs to the Conway normal form of the
\emph{exponent} $\gamma$. It is not coefficient extraction from an
unrelated field element $f$. Addition of surreal normal forms shows that
$q_\eta:\No\to\R$ is additive. If
$0<\delta\bb\omega^\eta$, then
\begin{equation}\label{eq:qvanish}
 q_\eta(C_\delta)=0.
\end{equation}
Indeed, if $0<|\gamma|\le n\delta$, the leading Conway exponent of
$\gamma$ is no larger than that of $\delta$. Since
$\delta\bb\omega^\eta$, this exponent is strictly below $\eta$,
so $q_\eta(\gamma)=0$. The case $\gamma=0$ is immediate.

\begin{example}[A concrete omnific-preserving nonmonomial image]
Take $\delta=1$ and $\theta=q_1$, the coefficient of $\omega$ in an
exponent. It kills every finite exponent and hence $C_1$.
The associated shear satisfies
\begin{equation}\label{eq:omegaexample}
 \sigma(\omega^\omega)
 =\omega^\omega\exp(-\omega^{-1})
 =\sum_{n\ge0}\frac{(-1)^n}{n!}\,\omega^{\omega-n}.
\end{equation}
Every exponent $\omega-n$ is positive, so the image is an omnific
integer. It is not a monomial. The automorphism fixes $\omega$, every
ordinary real, and every leading term. Formula \eqref{eq:omegaexample}
is a Hahn sum, not a claim of convergence in the fine topology of all
surreals.
\end{example}

\subsection{A parameter-avoidance lemma}

\begin{lemma}\label{lem:avoid}
Let $P\subseteq\No$ be a set. There is an ordinal $\eta>0$ such that
$q_\eta$ vanishes on every Hahn exponent occurring in every member of
$P$.
\end{lemma}
\begin{proof}
Let $E$ be the union of the Hahn supports of the members of $P$.
This is a set of surreal exponents. Let $T$ be the union of the Conway
normal-form exponent supports of the members of $E$. This is again a
set. Choose an ordinal $\eta>0$ strictly above every member of $T$;
for example an ordinal above their birthdays and zero suffices.
Then no exponent in $E$ has an $\omega^\eta$ coefficient.
This argument bounds a set of supports; it does not select a bound for
all surreal exponents at once.
\end{proof}

\begin{theorem}[Set-parameter invisibility]\label{thm:parameters}
For every set $P\subseteq\No$, there is a nonidentity strong
$1$-automorphism $\sigma$ of $\No$ such that
\[
 \sigma(\Oz)=\Oz,\qquad \sigma|_P=\id,
\]
and $\sigma$ sends a positive Conway monomial to a nonmonomial in
$\Oz$. Consequently the Conway monomial class is not definable in
$(\No,+,\cdot,<,\Oz)$ with parameters from any set $P$.
\end{theorem}
\begin{proof}
Choose $\eta$ as in \cref{lem:avoid}, put $H=\omega^\eta$, and take
$\delta=1$, $\theta=q_\eta$. Since $\eta>0$, $\theta(C_1)=0$.
The shear $\sigma=\sigma_{1,\theta}$ therefore preserves $\Oz$.
It fixes each monomial whose exponent is in $E$, so strongness implies
that it fixes every member of $P$.

But $\theta(-H)=-1$, and hence
\begin{equation}\label{eq:moveH}
 \sigma(t^{-H})=t^{-H}\exp(-t)
              =\sum_{n\ge0}\frac{(-1)^n}{n!}t^{-H+n}.
\end{equation}
All these exponents are negative because $H$ is positive infinite.
The nonzero terms at $-H$ and $-H+1$ already show that this is not a
monomial. Its positive leading coefficient makes it positive.

Any set definable with parameters from $P$ is preserved by every
automorphism fixing $P$ pointwise. The witness just constructed fails to
preserve the monomial class, proving nondefinability. A first-order
formula uses only finitely many parameters; the theorem proves the
stronger invariance failure even after an entire prescribed set is fixed.
\end{proof}

\subsection{No definable valuation representative of any kind}

There is a stronger conclusion than the failure to define the canonical
monomials. Write $x\sim_v y$ for $v(x)=v(y)$ on $\No^\times$,
equivalently $x/y\in\OO^\times$. Call a function
$s:\No^\times\to\No^\times$ a \emph{valuation representative rule} if
\begin{equation}\label{eq:sectionconditions}
 v(s(x))=v(x),\qquad x\sim_v y\Longrightarrow s(x)=s(y).
\end{equation}
For a set-sized field this is exactly a section of the quotient by
$\OO^\times$, written as a function on representatives. The formulation
\eqref{eq:sectionconditions} also makes sense on the full surreal class
without treating its proper-class equivalence classes as elements.
No multiplicativity, positivity, or monomial-image assumption is imposed.

\begin{lemma}[An entire valuation fiber has no fixed point]\label{lem:fiber}
For the shear in \cref{thm:parameters}, every $f$ with $v(f)=-H$ is
moved. More precisely,
\[
 v(\sigma f-f)=-H+1,
 \qquad\lc(\sigma f-f)=-\lc(f).
\]
\end{lemma}
\begin{proof}
The leading term of $f$ is $a t^{-H}$, with $a\ne0$.
Its first derivative under $D_{1,\theta}$ is $-a t^{-H+1}$.
All derivative contributions of later input terms have exponents larger
than $-H+1$, and every higher iterate of the leading term has exponent
$-H+n$ with $n\ge2$. Because $\theta(1)=0$, the iterate formula
\eqref{eq:iterate} applies. Thus the indicated first contribution cannot
cancel in $\Exp(D)f-f$.
\end{proof}

\begin{theorem}[No definable section]\label{thm:nosection}
For every set of parameters $P\subseteq\No$, no $P$-definable function
on $\No^\times$ satisfies \eqref{eq:sectionconditions}.
\end{theorem}
\begin{proof}
Suppose $s$ is such a function. Use the shear $\sigma$ and $H$ from
\cref{thm:parameters}, and let $x=t^{-H}$. Since $\sigma$ fixes all
valuations, $\sigma x\sim_v x$, and hence $s(\sigma x)=s(x)$.
Definability over fixed parameters gives
$s(\sigma x)=\sigma(s(x))$. Therefore $s(x)$ would be a fixed point
of $\sigma$ in the fiber of value $-H$, contrary to \cref{lem:fiber}.
\end{proof}

Theorems~\ref{thm:recovery} and \ref{thm:nosection} give a precise
boundary: the pair determines the finite/purely-infinite splitting and
ordinary constants, but does not determine a definable presentation of
all value classes by chosen elements. The nondefinability is relative to
the field-and-integer-part language. Conway's normal forms remain
set-theoretically defined in the ambient foundational construction.

\section{Fixed fields and arbitrarily large stabilizers}\label{sec:large}

\subsection{The exact common fixed field}

\begin{theorem}\label{thm:fixed}
The elements of $\No$ fixed by every strong omnific-preserving
$1$-automorphism are exactly the ordinary real numbers:
\[
 \Fix(\UU_\Oz)=\R.
\]
The same fixed-field statement holds if all field automorphisms preserving
$\Oz$ are allowed. The corresponding common fixed elements inside
$\Oz$ are exactly $\Z$.
\end{theorem}
\begin{proof}
All elements of $\UU_\Oz$ fix $\R$ by definition, and all unrestricted
pair automorphisms do so by \cref{cor:allautos}.
Let $f\notin\R$. Some nonzero $\gamma\in\supp(f)$ has a nonzero
Conway coefficient $q_\eta(\gamma)$. Put
\[
 \theta=q_\eta,\qquad \delta=\omega^{\eta-1}.
\]
Then $\delta>0$, $\theta(C_\delta)=0$, and $\theta(\delta)=0$.
The subset of $\supp(f)$ on which $\theta$ is nonzero has a least
member $\lambda$. By \eqref{eq:iterate}, the first nonzero contribution
to $(\Exp(D_{\delta,\theta})-\id)f$ is
\[
 a_\lambda\theta(\lambda)t^{\lambda+\delta}.
\]
Later affected input terms give larger exponents, and higher iterates
start at $\lambda+2\delta$. Thus this term does not cancel, and $f$
is moved by a member of $\UU_\Oz$. This proves the first equality.
The larger collection has the same fixed field because it contains
$\UU_\Oz$ and still fixes $\R$. Intersecting with $\Oz$ gives $\Z$.
\end{proof}

\subsection{Faithful additive product actions fixing parameters}

\begin{theorem}[An additive product of arbitrary set size]\label{thm:product}
For every set $P\subseteq\No$ and every set cardinal $\kappa$, there
is a faithful action of the additive group $\R^\kappa$ by strong
$1$-automorphisms preserving $\Oz$ and fixing $P$ pointwise. The group
is the full Cartesian product, not just the finite-support direct sum.
\end{theorem}
\begin{proof}
Regard $\kappa$ as its initial ordinal. The zero-cardinal case is
immediate. Let $E$ and $T$ be the support sets used in
\cref{lem:avoid}. Choose an ordinal $\beta$ strictly above $\kappa$
and every member of $T$. For $j<\kappa$ put
\[
 \delta_j=\omega^j,\qquad H_j=\omega^{\beta+j},
 \qquad\theta_j=q_{\beta+j},
\]
where $\beta+j$ is ordinary ordinal addition in this indexing formula.
These ordinals are distinct and at least $\beta$. The shifts $\delta_j$
form a strictly increasing, hence well-ordered, set. The functionals
satisfy
\begin{equation}\label{eq:orthogonal}
 \theta_j|_E=0,\quad \theta_j(C_{\delta_j})=0,\quad
 \theta_j(\delta_\ell)=0,\quad
 \theta_j(H_\ell)=\begin{cases}1&j=\ell,\\0&j\ne\ell.\end{cases}
\end{equation}
For $\mathbf a=(a_j)_{j<\kappa}\in\R^\kappa$, define
\begin{equation}\label{eq:Da}
 D_{\mathbf a}(t^\gamma)
  =t^\gamma\sum_{j<\kappa}a_j\theta_j(\gamma)t^{\delta_j},
\end{equation}
and extend strongly. Proposition~\ref{prop:uniform} proves that this is
well defined and contracting, and that its exponential preserves $\Oz$.
Every exponent in $E$ is annihilated, so $D_{\mathbf a}$ annihilates
$P$ and its exponential fixes $P$.

The identities $\theta_j(\delta_\ell)=0$ show, using the bracket law
or directly the product rule, that the $D_{\mathbf a}$ commute. Since
$D_{\mathbf a}+D_{\mathbf b}=D_{\mathbf a+\mathbf b}$, the map
$\mathbf a\mapsto\Exp(D_{\mathbf a})$ is a group homomorphism.
Moreover,
\begin{equation}\label{eq:faithful}
 \Exp(D_{\mathbf a})(t^{H_j})
       =t^{H_j}\exp(a_jt^{\delta_j}).
\end{equation}
If this action is the identity, the coefficient at $H_j+\delta_j$
in \eqref{eq:faithful} gives $a_j=0$ for every $j$. The action is faithful.

For arbitrary $\mathbf a$, all shifts still lie in the one set
$\{\delta_j:j<\kappa\}$. The support lemma supplies finiteness of each
coefficient across all word lengths. There is therefore no restriction
that only finitely many $a_j$ be nonzero. The result is a class relation
encoding a set-indexed action, not a set of proper-class graphs.
\end{proof}

\subsection{No finite solvability bound survives named parameters}

\begin{theorem}\label{thm:nonsolvable}
For every set $P\subseteq\No$, the pointwise stabilizer of $P$ inside
$\UU_\Oz$ is nonsolvable. More explicitly, for every ordinary $r\ge2$
it contains a subgroup of derived length exactly $r-1$.
\end{theorem}
\begin{proof}
Choose an ordinal $\beta$ above all exponents in the set $T$ from
\cref{lem:avoid}. Choose increasing ordinals
$\eta_1<\cdots<\eta_r$, all at least $\beta$, and set
\[
 e_j=\omega^{\eta_j},\qquad\theta_i=q_{\eta_i}.
\]
Then $e_1\bb\cdots\bb e_r$,
$\theta_i(e_j)$ is the Kronecker delta, and every $\theta_i$ kills
the parameter exponents $E$. For $j<i$, the derivation
$D_{e_j,\theta_i}$ preserves the boundary ring and annihilates $P$.
Take the subgroup generated by all
$\Exp(sD_{e_j,\theta_i})$ with $s\in\R$.

The bracket construction \eqref{eq:balanced} is unchanged, so
\cref{lem:mixed} proves that this subgroup has derived length at least
$r-1$. To see the upper bound, use the finite filtration by the shift
scales $C_{e_1},\ldots,C_{e_{r-1}}$ on the strongly closed Lie algebra
generated by these derivations. Sums of the positive shifts have the
largest scale of their summands; equal-scale brackets vanish. Thus the
proof of \cref{thm:rank} applies to these $r-1$ layers, even though the
ambient exponent group has many other convex subgroups. Its BCH group
has derived length at most $r-1$ and contains the stated subgroup.
All its generators fix $P$. Since $r$ is arbitrary, the full pointwise
stabilizer has no finite derived length.
\end{proof}

\subsection{The ring of omnific integers on its own}

\begin{proposition}\label{prop:fractions}
Every surreal is a quotient of two omnific integers. Every unital ring
automorphism of $\Oz$ extends uniquely to a field automorphism of
$\No$ preserving $\Oz$. Consequently the common fixed subring of all
ring automorphisms of $\Oz$ is $\Z$, and the Conway monomials lying
in $\Oz$ are not definable in the pure ring $\Oz$ with parameters from
any set.
\end{proposition}
\begin{proof}
For $f\in\No$, its Hahn support is a set, so choose $h>0$ above all
its exponents. Then both $t^{-h}$ and $t^{-h}f$ have strictly negative
support and lie in $\Oz$, giving the quotient representation.
An automorphism $\alpha$ of this integral domain extends by
$a/b\mapsto\alpha(a)/\alpha(b)$; cross multiplication proves
well-definedness, and the fraction-field property proves uniqueness and
surjectivity. The fixed-subring assertion follows from
\cref{thm:fixed}. For nondefinability, restrict the witness of
\cref{thm:parameters} to $\Oz$, with the parameter set contained in
that ring. It moves an omnific monomial to a nonmonomial.
\end{proof}
\status{The fraction-field property already occurs in the repository
\cite{RepoOz}. Its proof depends on the full surreal exponent class;
a fixed Hahn field can have an unbounded support and need not be the
fraction field of its nonpositive-support ring. The automorphism and
nondefinability consequences are the uses made here.}

\section{Surcomplex and Gaussian omnific consequences}\label{sec:complex}

Put
\[
 K_\C=\C((t^{\No}))\cong\No(i),\qquad
 A_\C=\Z[i]\oplus\C((t^{\No_{<0}}))=\Oz+i\Oz.
\]
The identification takes coefficientwise real and imaginary parts. The
union of the two real supports is well ordered, so it works in both
directions. Complex conjugation is coefficientwise, and $A_\C$ is the
ring of Gaussian omnific integers. It is not an ordered ring, and its
coefficient field is $\C$, not $\R$.

\begin{theorem}[Complex stabilizer and conjugation-compatible witnesses]
\label{thm:complex}
The stabilizer criterion of \cref{thm:stabilizer} applies to
$\C$-linear strong $1$-automorphisms preserving $A_\C$, with exactly
the same vanishing condition $c_\delta(C_\delta)=0$.
All real-coefficient shears and actions constructed in
\cref{sec:parameters,sec:large} extend to $\No(i)$, commute with
conjugation, and preserve $A_\C$.
For every set of surcomplex parameters they give both nondefinability
of the Conway monomials and failure of a definable valuation representative
rule, even in the language additionally naming the real axis and
conjugation. Their common fixed field, as $\C$-linear strong
$1$-automorphisms, is exactly $\C$.
\end{theorem}
\begin{proof}
The general stabilizer theorem was proved for every characteristic-zero
coefficient field and every unital coefficient subring, so apply it with
$(k,D)=(\C,\Z[i])$.
The functionals $q_\eta$ have real values. Thus their strong derivations
commute with coefficientwise complex conjugation, and so do their
exponentials. They preserve the real axis because their restriction
there is the original real automorphism.

For a set of complex parameters use the union of their real and imaginary
Hahn exponent supports in \cref{lem:avoid}. The resulting shear fixes
each parameter. The moved monomial \eqref{eq:moveH} remains real, so
the nondefinability proof is unchanged. In the valuation fiber of $-H$,
\cref{lem:fiber} holds with any nonzero complex leading coefficient;
its noncancellation argument does not require an order on coefficients.
The proof of \cref{thm:nosection} therefore applies verbatim to a
putative representative rule.

Finally every such $\C$-linear automorphism fixes $\C$, and for a
nonconstant complex series the proof of \cref{thm:fixed} uses a nonzero
exponent coefficient functional and a nonzero complex leading
coefficient. Their product is nonzero. The same real-coefficient shear
moves that series, proving the fixed-field assertion.
\end{proof}

For these real-axis-preserving maps, the surreal-valued modulus is
\emph{equivariant}:
\[
 |\sigma(z)|=\sigma(|z|).
\]
Indeed, $|z|$ is the unique nonnegative square root of
$z\overline z$, and $\sigma$ preserves multiplication, conjugation,
and order on $\No$. This is not a statement that every modulus is
fixed pointwise.

\subsection{Defining the complex constants from the Gaussian ring}

\begin{proposition}\label{prop:Cdef}
The ordinary complex field $\C$ is parameter-free definable in the
pure field-and-ring pair $(\No(i),A_\C)$. The common fixed field of
\emph{all} field automorphisms preserving $A_\C$ is $\Q$.
\end{proposition}
\begin{proof}
Let $B_\C=\C((t^{\No_{\le0}}))$ and
$\J_\C=\C((t^{\No_{<0}}))$.
If $u,v\in A_\C$ satisfy $u^2-2v^2=1$, the same factorization over
$B_\C$ and \cref{lem:units} give $u,v\in\Z[i]$.
Furthermore $u\ne0$: otherwise $v^2=-1/2$, impossible for a Gaussian
integer. Formula \eqref{eq:Jformula}, with $A_\C$ replacing $A$,
therefore defines $\J_\C$. One direction divides a purely infinite
series by the nonzero Gaussian integer $u$. In the converse, write the
constant coefficient as $a+bi$, $a,b\in\Z$, and choose an ordinary
positive Pell first coordinate $u>\max\{|a|,|b|\}$. Divisibility in
$\Z[i]$ by this ordinary integer forces $a=b=0$.
The idealizer proof now gives $B_\C$, and its units together with zero
are exactly $\C$.

The internal shears of \cref{thm:complex} show that a common fixed
point of all pair automorphisms must lie in $\C$. Every field
automorphism of $\C$ acts coefficientwise on $K_\C$ and preserves
$A_\C$: it sends $i$ to $i$ or $-i$, so preserves $\Z[i]$.
The common fixed field of these ordinary coefficient automorphisms is
$\Q$. For completeness, a nonrational algebraic number is moved by a
conjugate embedding extended to an algebraic closure and then to $\C$.
A transcendental number can be included in a transcendence basis over
$\Q$; translating that basis element by one and fixing the others
extends to an automorphism of $\C$. These standard algebraic-closure
extension arguments show that every complex number outside $\Q$ is
moved. All characteristic-zero field automorphisms fix $\Q$, proving
the claim.
\end{proof}

The real and complex fixed-field statements differ for a structural
reason. A real-field automorphism preserves order and cannot move an
ordinary real constant. A complex-field automorphism has no such
restriction unless additional real structure is required. Our explicit
nondefinability witnesses are stronger than needed in this respect:
they fix all of $\C$ and commute with conjugation.

\section{What the results settle, and what remains outside their scope}
\label{sec:scope}

The article answers three concrete structural questions. First, the
integer-part condition on a strong infinitesimal automorphism is exactly
a convex-scale annihilation condition on its logarithm. Second, the
strength of the residual symmetry is measured in finite rank by a sharp
solvability bound and in the full surreal class by arbitrarily large
parameter-fixing actions. Third, the omnific predicate recovers the
ordinary coefficient fields and zero-cut truncation but not a definable
valuation cross-section.

These conclusions impose several useful constraints on future approaches.
An attempt to reconstruct Conway monomials by a first-order formula in
the pair cannot succeed, even after choosing a set of additional
surreal constants. An attempt to classify strong $\Oz$-automorphisms
by coefficient and exponent actions alone misses the nontrivial
higher-rank factor $\UU_\Oz$. On the other hand, an alleged rank-one
nontrivial $1$-automorphism preserving the full divisible-exponent
integer part must violate one of the support hypotheses; the translation
example shows exactly how this happens.

Several different questions are deliberately not conflated with these
answers. We have not classified arbitrary nonstrong automorphisms of the
pair. The definability recovery does apply to them, but the logarithmic
classification requires strongness. We have not claimed that the shears
preserve the Gonshor exponential; the established triviality of the
strong exponential-field $1$-automorphism subgroup
\cite[Proposition 5.2]{KKS} is consistent with their being nontrivial
plain-field automorphisms. No new factorization theorem for omnific
integers is claimed; modern factorization work such as \cite{LM}
addresses a different problem. Finally, the arithmetic and model-theoretic
consequences proved here do not identify $\Oz$ with a model of full
ordinary integer arithmetic.

A natural next mathematical problem is to characterize which
\emph{nonstrong} automorphisms of $(\No,\Oz)$ commute with all
set-indexed normal-form sums. The present article supplies necessary
invariants and an exact answer once that additional property is imposed,
but it does not assert that the property follows from the first-order
pair alone.

\appendix
\section{Repository comparison and provenance ledger}\label{app:audit}

The repository was inspected on September 23, 2026, at the pinned revision
\begin{center}\small\pin.\end{center}
The inspected documentation includes the root README, the report
catalogue \path{docs/README.md}, the initial scope and results of
\path{docs/surreal/omnific-diophantine-geometry/article.tex}, and the
structural conventions and literature discussion of
\path{docs/surcomplex/surcomplex-field-automorphisms/article.tex}.
The catalogue also identified the related reports on set-sized omnific
quotients, exponential automorphism rigidity, and single-dilation Hahn
support. This was a targeted comparison, not a line-by-line audit of
all source manuscripts or Lean modules. The catalogue itself distinguishes
AI-assisted manuscripts, proof review, and implemented formalization.

\begin{longtable}{@{}>{\raggedright\arraybackslash}p{0.31\textwidth}>{\raggedright\arraybackslash}p{0.61\textwidth}@{}}
\toprule
Ingredient or result & Status in this article\\
\midrule
\endhead
Conway normal forms, $\Oz$, and Hahn support arithmetic
& Classical background; cited to Conway and Gonshor.\\[4pt]
Strong operator exponential--logarithm and BCH
& Imported from \cite{BKKPS}; surreal extension from \cite{KKS}.
No new proof of that general theorem is claimed.\\[4pt]
Ambient strong valued automorphism decomposition
& Established in \cite{KS}; its restriction to this stabilizer is
proved directly in \cref{thm:factorization}.\\[4pt]
Constant-product rigidity, Pell rigidity, and
$\operatorname{Frac}(\Oz)=\No$
& Already used or proved in \cite{RepoOz}; reproduced with attribution.\\[4pt]
Convex-scale logarithmic criterion
& Proposed additional theorem, proved in
\cref{thm:derivation,thm:stabilizer}.\\[4pt]
Finite-rank derived length and optimal examples
& Proposed additional results, with separate Lie and group proofs in
\cref{thm:rank,thm:sharp}.\\[4pt]
Pell-divisibility ideal definition and recovery of constants
& Additional definability deductions in
\cref{prop:Jdef,thm:recovery,prop:Cdef}.\\[4pt]
Set-parameter nondefinability and no representative rule
& Explicit additional constructions and proofs in
\cref{thm:parameters,thm:nosection}.\\[4pt]
Arbitrary product actions, nonsolvability, and common fixed fields
& Additional consequences with direct witnesses in
\cref{thm:fixed,thm:product,thm:nonsolvable,thm:complex}.\\[4pt]
Lean verification
& None supplied or claimed for these new statements. The finite Python
checks are not substitutes for formal proof.\\
\bottomrule
\end{longtable}

Primary-source checking used the 2025 version of \cite{BKKPS}, the
April 23, 2026 version of \cite{KKS}, and the cited Hahn-field,
truncation, and factorization literature. Searches for combinations of
omnific integers, automorphisms, integer parts, convex subgroups, and
definability did not locate the exact stabilizer-and-recovery package
proved here. Search coverage is not a mathematical guarantee of novelty;
independent literature review remains appropriate before a publication
priority claim.

\section{Finite verification and a proof-dependency audit}\label{app:checks}

\subsection{The included exact-arithmetic program}

The file \path{verify_finite_identities.py} uses only the Python standard
library. Its coefficients are exact rational numbers, not floating-point
approximations. It checks the homogeneous product rule and bracket formula
on finite Laurent polynomials, the balanced derived-series recursion,
finite-order shear exponential identities, the rank-one escape example,
and the Pell recurrence. It also expands balanced group commutators with
independent square-zero formal parameters. This last check verifies that
the mixed coefficient is the stated iterated Lie bracket for several
ranks, rather than inferring the group assertion from the Lie assertion.
Run it with
\begin{verbatim}
python3 verify_finite_identities.py
\end{verbatim}
A nonzero exit status or an assertion failure signals a failed check.
The delivered \path{verification_output.txt} records an actual run.

These tests are finite certificates for displayed identities. They do
not decide the well-ordering of arbitrary supports, establish the general
operator correspondence, verify a class-theoretic definition, or certify
historical novelty. The proofs of those parts are the written arguments
and their explicitly cited prerequisites.

\subsection{Points at which an apparently plausible proof can fail}

\paragraph{Dividing an exponent, not a coefficient.}
The implication from vanishing on $(-\delta,0)$ to vanishing on
$C_\delta$ uses divisibility of the \emph{exponent group}. The
characteristic-zero hypothesis separately permits rational scalar
coefficients. Neither assumption replaces the other.

\paragraph{Taking the logarithm in the right ring.}
The logarithm has coefficients $1/n$. We first pass from $A_D$ to the
$k$-algebra $B$, take the logarithm there, and only then deduce that its
nonconstant images lie strictly in $\J$. This avoids a false claim that
$A_D$ is closed under arbitrary rational multiplication.

\paragraph{Keeping a negative exponent negative at every iterate.}
It is not enough that $\gamma+\delta<0$ once. The annihilation condition
gives $-\gamma>n\delta$ for every ordinary $n$ whenever the functional
at $\gamma$ is nonzero. Thus the complete exponential, not only a finite
Taylor polynomial, stays on the correct side of the boundary.

\paragraph{Controlling full supports.}
The general classification starts with a strong operator. The explicit
many-shift constructions instead have one well-ordered positive shift
set, allowing the Neumann word lemma to prove every required
interchange. A family of additive functionals without either condition
would not be an adequate construction.

\paragraph{Not turning arbitrary Lie results into abstract group results.}
The upper bound uses BCH and strongly closed ideals. The lower bound uses
independent parameters, a nonzero mixed coefficient, and a finite
coefficient polynomial. These are distinct arguments, not a tacit
identification of abstract and closed derived subgroups.

\paragraph{No class of all class functions.}
Each automorphism and each action is encoded separately. The section
obstruction uses a function on nonzero elements satisfying
\eqref{eq:sectionconditions}, so no proper-class coset is treated as an
element of another class. The parameter-avoidance proof uses two unions
of set-sized supports, then chooses one ordinal bound.

\subsection{A possible formalization order}

A proof-assistant development can first treat a set-sized divisible ordered
abelian group. The minimal new algebraic ingredients are the additive
coefficient-function lemma, the principal-convex-subgroup vanishing lemma,
and the no-crossing support argument. Strong summability and an
operator-evaluation interface then yield the stabilizer theorem.
The finite-rank filtration and homogeneous bracket identities form a
separate algebraic layer. The mixed-parameter lemma can be formalized
using multivariate polynomials modulo the squares of their variables.

The definability argument should be treated independently: formalize the
Pell guard, the idealizer equality, the unit calculation, and the unique
negative/constant/positive decomposition. Only after the actual surreal
normal-form bridge is available should one instantiate exponent coefficient
functionals and the set-parameter bound. This is a proposed dependency
order, not a report that those modules have been implemented or compiled.

\clearpage
\section{Notation reference}\label{app:notation}
\begingroup
\small
\begin{longtable}{@{}>{\raggedright\arraybackslash}p{0.25\textwidth}>{\raggedright\arraybackslash}p{0.67\textwidth}@{}}
\toprule Symbol & Meaning\\\midrule\endhead
$\Gamma$ & A nonzero divisible ordered additive exponent group.\\[2pt]
$t^\gamma$ & Hahn monomial; on $\No$, it means $\omega^{-\gamma}$.\\[2pt]
$v(f),\lc(f)$ & Least Hahn exponent and its nonzero coefficient.\\[2pt]
$B,\J$ & Series supported at nonpositive and strictly negative exponents.\\[2pt]
$A_D$ & The ring $D\oplus\J$, where $D$ is a unital subring of $k$.\\[2pt]
$\OO,\mm$ & The nonnegative- and positive-support valuation pieces.\\[2pt]
$C_\delta$ & The principal convex subgroup of $\Gamma$ generated by
$\delta>0$.\\[2pt]
$\delta\bb\epsilon$ & $n\delta<\epsilon$ for every positive ordinary
integer $n$.\\[2pt]
$c_\delta(\gamma)$ & Coefficient of $t^{\gamma+\delta}$ in
$\partial(t^\gamma)$.\\[2pt]
$D_{\delta,\theta}$ & The homogeneous derivation
$t^\gamma\mapsto\theta(\gamma)t^{\gamma+\delta}$.\\[2pt]
$\LL_A,\UU_{A_D}$ & Boundary-preserving contracting strong derivations,
and their strong $1$-automorphism exponentials.\\[2pt]
$P_-,P_0,P_+$ & Projection to negative support, constant coefficient,
and positive support.\\[2pt]
$q_\eta(\gamma)$ & The coefficient of $\omega^\eta$ in the Conway
normal form of the exponent $\gamma\in\No$.\\[2pt]
$\Exp,\Log$ & Operator exponential and logarithm.\\[2pt]
$\exp(u)$, $v(u)>0$ & The strongly summable Hahn series $\sum u^n/n!$.\\[2pt]
$A_\C$ & $\Oz+i\Oz$, the Gaussian omnific integers.\\
\bottomrule
\end{longtable}
\endgroup

\clearpage
\phantomsection
\addcontentsline{toc}{section}{References}
\begin{thebibliography}{99}

\bibitem{Conway}
J.~H.~Conway,
\emph{On Numbers and Games},
L.M.S. Monographs, Academic Press, 1976.

\bibitem{Gonshor}
H.~Gonshor,
\emph{An Introduction to the Theory of Surreal Numbers},
London Mathematical Society Lecture Note Series 110,
Cambridge University Press, 1986.

\bibitem{KS}
S.~Kuhlmann and M.~Serra,
\emph{The automorphism group of a valued field of generalised formal power series},
Journal of Algebra \textbf{605} (2022), 339--376.
\href{https://doi.org/10.1016/j.jalgebra.2022.04.023}{doi:10.1016/j.jalgebra.2022.04.023}.
Preprint \href{https://arxiv.org/abs/2107.03362v3}{arXiv:2107.03362v3}.

\bibitem{BKKPS}
V.~Bagayoko, L.~S.~Krapp, S.~Kuhlmann, D.~Panazzolo, and M.~Serra,
\emph{Automorphisms and derivations on algebras endowed with formal infinite sums},
\href{https://arxiv.org/abs/2403.05827v2}{arXiv:2403.05827v2},
September 22, 2025. Especially Theorems 3.11 and 3.13.

\bibitem{KKS}
E.~Kaplan, L.~S.~Krapp, and M.~Serra,
\emph{Decomposing the automorphism group of the surreal numbers},
\href{https://arxiv.org/abs/2509.22374v3}{arXiv:2509.22374v3},
April 23, 2026. Especially Fact 4.2 and Proposition 5.2.

\bibitem{Camacho}
S.~Camacho,
\emph{Truncation in Hahn Fields is Undecidable and Wild},
\href{https://arxiv.org/abs/1706.03722}{arXiv:1706.03722}, 2017.

\bibitem{LM}
S.~L'Innocente and V.~Mantova,
\emph{A factorisation theory for generalised power series and omnific integers},
\href{https://arxiv.org/abs/1710.07304v5}{arXiv:1710.07304v5},
January 22, 2024.
\href{https://doi.org/10.1016/j.aim.2024.109513}{doi:10.1016/j.aim.2024.109513}.

\bibitem{Repo}
VladimirReshetnikov/Surreal,
\emph{Repository and research-report catalogue},
revision \pin, inspected September 23, 2026.
\href{https://github.com/VladimirReshetnikov/Surreal/tree/fb5c4b530e3ec04b5661347d51d5382526d5aa02}{Pinned repository}.

\bibitem{RepoOz}
Surreal project,
\emph{Omnific Integers and Omnific--Diophantine Geometry},
AI-assisted research manuscript, in \cite{Repo}.
\href{https://github.com/VladimirReshetnikov/Surreal/blob/fb5c4b530e3ec04b5661347d51d5382526d5aa02/docs/surreal/omnific-diophantine-geometry/article.tex}{Pinned manuscript source}.
Used for the prior constant-product, Pell-rigidity, and fraction-field results.

\bibitem{RepoAut}
Surreal project,
\emph{Surcomplex Field Automorphisms},
AI-assisted research manuscript, in \cite{Repo}.
\href{https://github.com/VladimirReshetnikov/Surreal/blob/fb5c4b530e3ec04b5661347d51d5382526d5aa02/docs/surcomplex/surcomplex-field-automorphisms/article.tex}{Pinned manuscript source}.
Used for comparison of structures preserved and class-map conventions.

\end{thebibliography}
\end{document}
